% Teil II: durchgerechneter Weg, Teile a--b.

\subsection*{Gegeben und gesucht}
Masseloser Balken, bei der Mitte auf einer Schneide drehbar gelagert; am
rechten Ende (Armlänge $a$) hängt $m_2$, am linken Ende ist ein Stück
der Länge $b$ senkrecht zum Balken abgeknickt, an dessen Ende $m_1$
hängt; Balken um $\alpha$ gegen die Horizontale geneigt.
Spezialfall: $a=2b$, $m_2=1,3\,m_1$. Gesucht: (a) Kräfte im Balken,
(b) Gleichgewichtsbedingung und der Winkel $\alpha$.

\subsection*{Ansatz}
Momentengleichgewicht um die Schneide (Drehpunkt), dort fällt die
Lagerkraft heraus. Hebelarme horizontal ansetzen.

\subsection*{Rechnung}
\textbf{a) Kräfte.} Auf den Balken wirken: die Gewichtskraft $m_2g$
vertikal nach unten am rechten Ende, $m_1g$ vertikal nach unten am Ende
des abgeknickten Stücks, und die Lagerkraft unter der Schneide (im
Drehpunkt). Balkenmasse und -dicke vernachlässigen.

\textbf{b) Gleichgewicht und Winkel.} Um die Schneide (Momentensumme
null). Horizontale Hebelarme: rechts $\ell_2=a\cos\alpha$; links beträgt
der Arm $a\cos\alpha$ (bis zum Ende) plus $b\sin\alpha$ (zusätzlich durch
das abgeknickte, nach aussen gerichtete Stück), also
$\ell_1=a\cos\alpha+b\sin\alpha$. Mit dem Drehsinn (rechts abwaerts /
links aufwaerts negativ) gibt
\begin{equation}
  -m_2g\,a\cos\alpha + m_1g\,(a\cos\alpha+b\sin\alpha)=0
  \;\Rightarrow\;
  m_2\,a\cos\alpha = m_1\,a\cos\alpha + m_1\,b\sin\alpha
\end{equation}
\begin{equation}
  (m_2-m_1)\,a\cos\alpha=m_1\,b\sin\alpha
  \;\Rightarrow\;
  \boxed{\;\tan\alpha=\frac{(m_2-m_1)\,a}{m_1\,b}\;}.
\end{equation}
Einsetzen von $a=2b$, $m_2=1,3\,m_1$:
\begin{equation}
  \tan\alpha=\frac{(1,3m_1-m_1)\cdot2b}{m_1\,b}=0,6
  \;\Rightarrow\;
  \boxed{\;\alpha=\arctan(0{,}6)\approx\qty{31.0}{\degree}\;}.
\end{equation}

\begin{ergebnis}
Gleichgewichtsbedingung $\sum M=0$ um das Lager;
$\tan\alpha=(m_2-m_1)/(m_1)\cdot a/b$. Für $a=2b$, $m_2=1,3m_1$:
$\tan\alpha=0{,}6$, also $\alpha\approx\qty{31.0}{\degree}$.
\end{ergebnis}

\begin{kontrolle}
\textbf{Grenzfall:} Sinkt der Abstand $m_2=m_1$ auf gleich grosse Massen,
wird $\tan\alpha=0$ und $\alpha=0$ (waagrechter Balken) -- symmetrische
Anordnung, plausibel. \checkmark

\textbf{Einheiten:} $\tan\alpha$ ist dimensionslos; $a$ und $b$ in m
kürzen sich unter dem Bruch. \checkmark
\end{kontrolle}

\begin{merksatz}
Im Drehmoment-Gleichgewicht um ein Lager verschwinden die Momente aller
Gewichte; horizontale Hebelarme sind $\ell\cos\alpha \pm$ ($b\sin\alpha$
für abgeknickte Teile). Der Winkel folgt aus
$\tan\alpha=(m_2-m_1)a/(m_1b)$.
\end{merksatz}
